\section{Background}\label{sec:background}

The problem of mandatory exception escalation in autonomous multi-agent systems draws from four distinct intellectual traditions: accountability architecture in distributed AI systems, human-in-the-loop design, fail-safe engineering in safety-critical systems, and the BX3 three-layer accountability model as the structural substrate that makes compulsory escalation architecturally enforceable. Each tradition offers essential conceptual resources, but each also contains a specific structural gap that the Bailout Protocol is designed to address at its root.

% ---------------------------------------------------------------
\subsection{Accountability in Multi-Agent AI Systems}\label{sec:accountability}

Accountability in multi-agent AI systems is the property that every consequential action can be attributed to a named principal, and that the chain of delegation connecting any action to its originating authority is traceable, verifiable, and unbreakable. The accountability problem in single-agent systems is straightforward by comparison: a single actor is responsible for all of the system's outputs. In multi-agent systems, where a task may be decomposed and delegated recursively across multiple spawned nodes, accountability becomes a structural challenge rather than merely a governance one. The question is not simply whether the system has policies requiring attribution --- it is whether the system's architecture guarantees that the attribution chain cannot be severed by any machine actor operating within it.

The literature on AI accountability has focused primarily on audit trails, provenance records, and governance frameworks as post-hoc solutions to this problem. ChainScore Labs argues that ``only cryptographically verifiable, on-chain provenance for training data, model weights, and inference outputs can create the immutable audit trail required for safety, compliance, and trust in autonomous AI systems'' \cite{on_chain_provenance_ai}. This blockchain-provenance approach achieves impressive properties: immutable timestamps, independently verifiable continuity, and cryptographic integrity of the record. However, its fundamental limitation is that it assumes the system being audited is the system that generated the record. An agent that can rewrite its own parent reference can present a provenance chain to auditors that diverges from the chain under which it actually operated --- and the audit will confirm the presented chain as valid, because the audit verifies the record's integrity, not the topology that produced it.

FlexBlok proposes blockchain-based distributed ledgers as independently verifiable, tamper-evident records for AI governance \cite{blockchain_ai_governance}. Like ChainScore, this approach provides strong properties for record integrity but inherits the same structural vulnerability: the ledger records events that the system chooses to emit. If a machine actor omits an event, rewrites a parent reference before emitting the next event, or presents a modified delegation chain to downstream verification services, the ledger confirms consistency with what was emitted --- not fidelity to what actually occurred.

The Human Delegation Provenance (HDP) protocol represents the most directly relevant prior work \cite{hdp2025arxiv}. HDP proposes that human authorization decisions travel through the delegation chain as first-class cryptographic artifacts: every HDP token encodes scope, duration, and constraints, and verification is performed at each chain traversal. The protocol's core contribution is establishing human authorization as a per-decision artifact rather than a one-time grant at chain initialization. However, HDP assumes a trustworthy delegation chain topology; it verifies the integrity of tokens traversing the chain but not the integrity of the chain itself. In a mutable-chain environment, a misbehaving agent can present a false ancestor chain before HDP token verification, and the token will verify correctly against the false topology because the token's own cryptographic integrity is independent of whether the chain topology is authentic. HDP closes the per-decision authorization gap but not the chain-integrity gap.

The structural conclusion from this body of work is that accountability requires a chain substrate that is itself immutable --- not merely a record that is append-only after the fact. An immutable parent chain means that the delegation topology is fixed at node creation and cannot be modified by any actor for the node's lifetime. The provenance record is then trustworthy because the topology it records is the topology that actually existed. The BX3 Framework's forensic ledger achieves the same independence properties as blockchain-based approaches --- append-only, hash-chained, independently recomputable --- but within the framework's own execution environment \cite{bx3framework}. This eliminates the operational overhead of external distributed consensus while providing equivalent auditability guarantees.

% ---------------------------------------------------------------
\subsection{Human-in-the-Loop Design}\label{sec:hitl}

Human-in-the-loop (HITL) design is the principle that human judgment must be available, required, or determinative at one or more points in an autonomous system's operation. The principle originated in safety-critical systems engineering --- aviation, nuclear process control, medical devices --- where the consequences of autonomous machine decision-making under failure conditions are potentially catastrophic and human oversight is the only credible backstop against unintended outcomes \cite{bansal2019}. Wiener's foundational analysis of cybernetics established the theoretical basis for this intuition: a control system that severs its feedback connection to a responsible human operator becomes, in Wiener's term, ``outside the circle of ethical accountability'' \cite{wiener1948}. The ethical boundary of autonomous operation, in other words, is defined by the retention of a communication channel to a human who can intervene.

In AI governance, HITL has been widely adopted as a design norm. The EU AI Act classifies high-risk AI systems as requiring ``meaningful human oversight'' \cite{eua2024}. NIST's AI Risk Management Framework recommends human role definition as a core function of AI governance \cite{nist2023}. IEEE standards for autonomous systems specify human override capability as a baseline requirement \cite{ieee2024autonomous}. These frameworks share a common structure: they define HITL as a policy requirement that system designers implement, and they rely on compliance assessment to verify implementation.

The structural gap in policy-based HITL is the distinction between \emph{discretionary} and \emph{compulsory} human oversight. In a discretionary model --- which describes essentially all current production frameworks --- the system or its operators decide whether to request human input, and the system may proceed autonomously if human input is unavailable, delayed, or not sought. This creates a failure mode that the accountability literature has not fully characterized: a system operating in a discretionary HITL mode is structurally identical to a system with no HITL requirement, from the perspective of the accountability chain. The human is present as a possible input but not as a compulsory terminus.

Trist and Bamforth's analysis of organizational change establishes that technological systems and social systems are co-constituted: introducing a new technological structure changes the responsibilities, authority, and accountability relationships of the humans embedded within it \cite{tristr81}. Applied to HITL design, this implies that if an autonomous system introduces a human as a discretionary consultant rather than as a compulsory decision terminus, the human's authority is structurally diminished by the architecture itself, regardless of what the system's policy documents claim. The human is present but not determinative.

Bansal et al. demonstrate in the context of autonomous vehicle safety that even well-designed HITL systems can fail in ways that are difficult to anticipate: human operators develop automated complacency, relying on the system to catch errors that the human should be catching, so that when the system fails in an unanticipated way the human is not prepared to intervene effectively \cite{bansal2019}. This finding has been replicated across aviation, medical, and industrial settings. The structural lesson is that HITL designed as a discretionary layer on top of autonomous operation does not reliably produce human oversight in practice, because human vigilance is not maintained when humans believe the system is handling the difficult cases.

The Bailout Protocol addresses this gap by making human oversight compulsory and structurally enforced rather than discretionary and policy-dependent. Entry to the protocol is not a decision the system makes --- it is a state the system enters when its operating conditions exceed its provisioned scope, independent of any assessment about whether human oversight is warranted. The human principal $h(n)$ is not one possible terminus among several; it is the only permissible terminus for every unresolved exception. This transforms HITL from a behavioral property of the system (humans may be consulted under specified conditions) into a structural property of the architecture (humans must be reached for every exception that exceeds machine authority).

% ---------------------------------------------------------------
\subsection{Fail-Safe Design in Safety-Critical Systems}\label{sec:failsafe}

Fail-safe design is the engineering principle that a system must transition to a safe state when it encounters a condition it cannot continue to handle correctly. The concept originates in safety-critical systems engineering, where the cost of failure is measured in human lives rather than economic loss. A fail-safe system does not merely tolerate failure --- it ensures that when failure exceeds the system's recovery capacity, the system degrades to a state that is safe by definition, not a state that is merely better than the worst alternative.

Dijkstra's seminal work on self-stabilizing distributed systems established the formal foundation for fail-safe reasoning in systems with multiple concurrent components \cite{dijkstra1982}. A self-stabilizing system has the property that, starting from any arbitrary initial state --- including states produced by arbitrary transient faults, arbitrary sequence of component failures, or deliberate adversarial manipulation --- the system eventually converges to a legitimate state within bounded time, and thereafter remains in legitimate states. The critical implication is that correctness is defined not only for the system's intended operational state but for every possible state the system might enter, including states produced by failure.

Fail-safe design goes beyond self-stabilization by adding a directional requirement: when a system cannot continue correctly, it must not continue operating in any capacity that could produce harm. In mechanical fail-safe systems, this is often achieved by returning to a passive state --- a valve closes, a brake engages, a circuit opens. In computational systems, the analogous principle is that the system must not continue autonomous operation when its operating conditions are undefined: it must halt, escalate, or degrade in a manner that is safe by construction.

The distinction between fault tolerance and fail-safe operation is architecturally significant. Fault tolerance assumes that failures are within known bounds and that the system can recover through redundant or alternative pathways. Fail-safe assumes that failures may exceed the system's recovery capacity and that the only defensible response is to surface the failure to a responsible human decision-maker. In safety-critical systems, this distinction is settled: when the bounds of safe autonomous operation cannot be established with confidence, fail-safe prevails.

The Bailout Protocol embodies fail-safe principles in a multi-agent AI architecture. When a node's Bounds Engine encounters a condition outside its provisioned operating range $R(n)$ that it cannot resolve within delegated authority, the node does not attempt autonomous recovery, does not continue operation in a degraded mode, and does not pass the exception to an intermediate machine actor for resolution. The node enters a defined safe state --- the ESCALATING state --- and the exception propagates to $h(n)$. This satisfies Dijkstra's self-stabilization criterion in a strengthened form: from any arbitrary exception state, every node necessarily reaches the designated safe state $h(n)$ within bounded time, and $h(n)$ is safe by definition because it involves human judgment \cite{dijkstra1982}. The Bailout Protocol additionally satisfies the fail-safe directional requirement: there is no intermediate state in which the node operates autonomously in an undefined or unsafe condition. Entry to ESCALATING is atomic and irreversible by any machine actor, and the node cannot re-enter autonomous operation without a binding resolution directive from $h(n)$.

% ---------------------------------------------------------------
\subsection{The BX3 Three-Layer Model}\label{sec:three-layer}

The \bxthree{} Framework is an accountability architecture for multi-agent AI systems that organizes every node into three immutable functional layers: the \purpose{} layer, the \bounds{} layer, and the \fact{} layer \cite{bx3framework}. This three-layer architecture is the structural substrate that makes the Bailout Protocol architecturally enforceable rather than merely policy-dependent. The functional separation between layers is not a design convenience; it is a structural necessity that produces the accountability guarantees without which the Bailout Protocol could not function.

The \purpose{} layer carries intent, judgment, and final authority over any decision that exceeds a node's provisioned scope. It defines the operating range $R(n)$ for each node, establishes the human accountability anchor $h(n)$, and receives all escalated exceptions. The \purpose{} layer is the exclusive source of intent in the node; no other layer may generate a directive that overrides it.

The \bounds{} layer carries bounded reasoning and constrained execution within $R(n)$. It evaluates whether any proposed action falls within the node's current authority, detects exception conditions that exceed $R(n)$, and initiates the Bailout Protocol when it cannot produce a \fact-layer-approved resolution within the bounded timeout $T_{\text{bounds}}$. The \bounds{} layer carries no judgment authority over exceptions --- it detects and signals but does not adjudicate.

The \fact{} layer carries deterministic enforcement and forensic auditability. It verifies proposed actions against the authorization ledger, enforces the immutable parent pointer via pre-commit hooks at node creation, executes the state transition relation of the Bailout Protocol's state machine, and maintains the append-only forensic ledger recording every protocol event as an immutable architectural record \cite{bx3framework}. The \fact{} layer carries no intent authority.

The functional separation between \purpose{}, \bounds{}, and \fact{} layers creates the structural condition that makes the Bailout Protocol not merely possible but compulsory. When an exception exceeds the \bounds{} layer's provisioned scope, no machine actor exists that has the authority to adjudicate it. The \bounds{} layer detects but carries no judgment authority; the \fact{} layer enforces but carries no intent authority. The decision must return to the \purpose{} layer, which carries the intent and final accountability. This is not a policy choice --- it is a structural consequence of the layer separation.

The three-layer model is minimal with respect to the Bailout Protocol's correctness guarantees. Removing the \purpose{} layer eliminates the exclusive terminus of escalation: machine actors could absorb exceptions and issue terminal directives, and the accountability chain would terminate in autonomous resolution rather than human judgment. Removing the \bounds{} layer converts the escalation trigger from a structural enforcement into a policy parameter that operational urgency can exceed. Removing the \fact{} layer makes the immutable parent pointer unenforceable and the forensic ledger unauditable. No proper subset of the three layers produces the same guarantees \cite{bx3framework}. The Bailout Protocol is therefore not a feature layered onto an existing architecture; it is co-constituted with the three-layer model itself --- each is the other's structural prerequisite.

The forensic ledger maintained by the \fact{} layer provides the append-only, hash-chained record of all bailout events, escalation events, and resolution directives that makes the Bailout Protocol independently verifiable. Every state transition, every timeout event, and every directive issued by $h(n)$ is recorded with cryptographic continuity to its predecessor. An external auditor can recompute the hash chain and confirm continuity without relying on the system's own operators or accepting their testimony. This is the property that blockchain-provenance systems seek to provide through external distributed consensus \cite{blockchain_ai_governance} \cite{on_chain_provenance_ai}, achieved within the framework's own execution environment with lower operational overhead and no dependency on consensus mechanisms outside the framework's own structural guarantees.
